\subsection{A prime-to-three cover resolving the small genus-two test}

\begin{result}{B14}{Proved}{The two ramified points in the genus-two test are separated}
Let $C/\mathbb F_3$ be the smooth projective model of
\[
                 y^3-y=x+x^{-1},
\]
and let $P,Q$ be its rational points over $x=0,\infty$, respectively.
Their point sections in the geometrically prime-to-three fundamental
group are not geometrically conjugate. This remains true after every
finite extension of $\mathbb F_3$. In particular their decomposition
groups are not conjugate-commensurate. We construct a finite \emph{\'etale}
cover of $C_{\mathbb F_{27}}$, with geometric monodromy of order prime
to three, having a rational point over $P$ and none over $Q$.
Thus this genus-two test is decided positively, although the
general higher-genus assertion retained in B9 is not a consequence
of this computation.
\par\emph{Passage:} cote 149, \S1, manuscript p.~17, scans 35--36,
the proposed recovery of points by prime-to-characteristic covers (F7).
\end{result}

\begin{proof}
We make the covers and the section-lifting comparison explicit.

\emph{The quadratic cover and its two quotients.}
Put $k=\mathbb F_{27}$ and choose $a\in k$ with $a^3-a=1$;
$T^3-T-1$ is irreducible over $\mathbb F_3$ since it has no root there.
Set $u=y-a$ and $z=x-x^{-1}$. The original equation gives
\[
 z^2=(u^3-u+1)^2-1
    =u(u-1)(u+1)(u^3-u-1).
 \tag{B14.1}
\]
Conversely $y=u+a$ and $x=(u^3-u+1+z)/2$, so this is an
isomorphism of function fields and hence of the smooth projective
models. The sextic has six distinct roots; this also checks $g(C)=2$.
As in B9, $\operatorname{div}(x)=3P-3Q$, and $[Q-P]$ has exact
order three in $J_C(\overline k)$: order one would give a degree-one
map from this genus-two curve to $\mathbb P^1$.

Let $D$ be the smooth projective curve with function field
\[
 k(u,v,w),\qquad
 v^2=u(u-1),\qquad
 w^2=G(u):=(u+1)(u^3-u-1),
 \qquad z=vw.
 \tag{B14.2}
\]
The last equality gives a degree-two morphism $f:D\to C_k$.
It is connected: neither $u(u-1)$ nor its quotient by the sextic
in (B14.1) is a square in $\overline k(u)$, as their simple zero
or pole valuations show. The divisor of $u(u-1)$ on $C_k$ is even.
Indeed its two zeros are hyperelliptic branch points and have order
two, while at each of $P,Q$ it has a pole of order two. All other
orders are zero. Quadratic Kummer theory in residue characteristic
three therefore proves that $f$ is \emph{\'etale} everywhere.

There is also a degree-two map $g:D\to E$, where
\[
             E:\quad w^2=G(u)=u^4+u^3+2u^2+u+2.
 \tag{B14.3}
\]
The quartic is square-free, so its projective model has genus one.
It has two distinct rational points $e_+,e_-$ at infinity,
labelled by $w/u^2=+1,-1$. The involutions of $D$ over $C$ and
$E$ are, respectively,
\[
 \tau(u,v,w)=(u,-v,-w),\qquad
 \sigma(u,v,w)=(u,-v,w).
\]
Their product $h(u,v,w)=(u,v,-w)$ has the conic
$B:v^2=u(u-1)$ as quotient. This conic is rational, with
\[
 t=v/u,\qquad u=(1-t^2)^{-1},\qquad v=t(1-t^2)^{-1}.
\]
Writing $W=w(1-t^2)^2$ gives the particularly simple model
\[
                D:\quad W^2=1-t^2-t^4-t^8.
 \tag{B14.4}
\]
The degree-eight polynomial is square-free. Hence $D\to B\simeq
\mathbb P^1$ is branched at eight points and $D$ is hyperelliptic
of genus three. This agrees with the unramified formula
$2g(D)-2=2(2g(C)-2)$.

\emph{The rational lifts and their elliptic images.}
At $P$ and $Q$, respectively, $z/u^3$ tends to $-1$ and $+1$.
Since in (B14.4) this quotient is $tW$, we can choose the rational
lift and list the two rational lifts of $Q$ as
\[
       P'=(t,W)=(1,-1),\qquad
       Q_+=(1,1),\qquad Q_-=(-1,-1).
 \tag{B14.5}
\]
These are all defined over $\mathbb F_3$ on the displayed model of
$D$; their map to the original $C$ is defined over $k$ because it
uses $a$. The formula $w/u^2=W$ gives
\[
            g(P')=g(Q_-)=e_-,\qquad g(Q_+)=e_+.
\]
The elliptic curve (B14.3) has exactly five $\mathbb F_3$-points:
above $u=0$ there are none, above $u=1$ there are two, above $u=-1$
there is one, and there are its two points at infinity. Therefore,
with $e_-$ as origin, $e_+-e_-$ has exact order five. It follows
that $[Q_+-P']\in J_D(\overline k)$ is not three-primary.

\emph{The other lift has Jacobian difference of exact order six.}
Put $\delta=[Q_--P']$. The two norm maps give
\[
        f_*\delta=[Q-P],\qquad g_*\delta=0.
\]
We verify the necessary assertion about their joint kernel instead
of assuming a decomposition of the Jacobian. On divisor classes,
pullback followed by norm satisfies
\[
             f^*f_*=1+\tau_*,\qquad
             g^*g_*=1+\sigma_*.
\]
The hyperelliptic involution $h_*$ is $-1$ on $J_D$: every divisor
$R+hR$ is a fibre of the degree-two map to $B$, so applying $h$
to a degree-zero divisor class negates it. Since $\tau\sigma=h$,
we have $\tau_*=-\sigma_*$, and consequently
\[
                  f^*f_*+g^*g_*=[2].
 \tag{B14.6}
\]
Thus the joint norm kernel is killed by two. Applying (B14.6) to
$3\delta$ proves $6\delta=0$, while the exact order of $f_*\delta$
shows that the order of $\delta$ is divisible by three.

It cannot have order three. Such an order would supply a function
$b$ on $D_{\overline k}$ with
$\operatorname{div}(b)=3Q_--3P'$, hence a degree-three morphism
to $\mathbb P^1$. Together with the degree-two hyperelliptic map,
this would give a birational map onto an integral curve of bidegree
$(2,3)$ in $\mathbb P^1\times\mathbb P^1$. To check birationality,
the degree of the function-field extension over the compositum of
the two rational subfields divides both two and three, so is one.
An integral curve of bidegree $(2,3)$ has arithmetic genus
$(2-1)(3-1)=2$, by adjunction; its normalization has genus at most
two. This contradicts $g(D)=3$. This argument uses function-field
degrees, so does not silently assume that the degree-three map is
separable. We conclude that $\delta$ has exact order six. Both
lifts in (B14.5) therefore have non-three-primary difference from
$P'$, over $k$ and over every algebraic extension of $k$.

\emph{The actual cover separating the fibres.}
Here is a finite cover constructed from the preceding classes.
Write $J=J_D$ over $k$, let $F$ be its $27$-power Frobenius,
and let $A=J(k)[3^\infty]$, viewed as a constant finite subgroup
scheme. The Lang isogeny $F-1:J\to J$ is finite \emph{\'etale}
with kernel $J(k)$: its derivative is $-1$, and its geometric
kernel is precisely the Frobenius-fixed finite group. Quotienting
its source by $A$ gives a finite \emph{\'etale} isogeny
\[
           \lambda:J/A\longrightarrow J,
           \qquad \ker\lambda=J(k)/A,
 \tag{B14.7}
\]
of degree prime to three. The comparison of rational images is
exact:
\[
                     \lambda((J/A)(k))=A\subset J(k).
 \tag{B14.8}
\]
Indeed, for a rational class $\bar b\in(J/A)(k)$, any geometric
lift $b\in J(\overline k)$ satisfies $F(b)-b\in A$, so its image
under $\lambda$ belongs to $A$. Conversely for $c\in A$ choose
$b$ with $F(b)-b=c$, using surjectivity of $F-1$; its class modulo
$A$ is fixed by Frobenius and is rational. This proves (B14.8).

Let $\alpha:D\to J$ be $R\mapsto[R-P']$. Then
\[
             V=D\times_{\alpha,J,\lambda}(J/A)
               \longrightarrow D\xrightarrow{f}C_k
 \tag{B14.9}
\]
is the promised cover. It has a rational point above $P'$ supplied
by $0\in J/A$. By (B14.8) and the proved orders, it has no rational
point over either $Q_+$ or $Q_-$. These are all points of $D$
above $Q$, so it has none over $Q$. Its geometric monodromy is
prime to three: the first map is a torsor under the finite group
$J(k)/A$, and the second is a Galois double cover. The Galois
closure of their composite embeds, over the double cover, into the
product of the first torsor group and its deck conjugate; its
geometric group is an extension of a subgroup of that product by
a group of order two. Every order involved is prime to three.
Taking a connected component containing the displayed rational
point is also permissible and preserves these properties.

For numerical reproducibility, the exact script
\texttt{code/F7\_test.py} enumerates
$\#D(\mathbb F_3)=6$, $\#D(\mathbb F_9)=10$,
$\#D(\mathbb F_{27})=36$. Newton identities and the functional
equation give its Frobenius polynomial
\[
 Z^6+2Z^5+2Z^4+4Z^3+6Z^2+18Z+27.
\]
Using the curve trace formula and the Jacobian point formula
\cite[III, Theorem 11.1; II, Theorem 1.1]{AB_MilneAV} yields
$\#J(\mathbb F_{27})=25920=2^6\cdot3^4\cdot5$.
Thus (B14.7) has degree $320$ and (B14.9) has degree $640$
before choosing a component. The proof of separation above does
not depend on these numerical counts beyond the elementary five
points on $E$.

\emph{Comparison with point sections and constant extensions.}
For any finite \emph{\'etale} cover whose geometric monodromy is
prime to three, the rational points of its fibre over $R\in C(k)$
are exactly the fixed points of the image of $s_R(G_k)$ in the
corresponding permutation representation. Geometrically conjugate
sections have isomorphic such $G_k$-sets. The different rational
fibre counts in (B14.9) therefore prove the required nonconjugacy
over $k$, and hence over $\mathbb F_3$.

To obtain the asserted persistence after every finite extension,
one may repeat the Lang construction with that finite field's
Frobenius, after adjoining $a$. The two non-three-primary divisor
classes retain their orders, so the same proof works. Equivalently,
the precise covering argument of B8 applies to $D\to C$: after
making equal sections literally equal by a geometric conjugator,
the fixed coset supplied by the rational lift $P'$ supplies a
rational lift $Q'$ with the same section in $\pi_1(D)$. This
is the full geometrically prime-to-three group of $D$, since
intersecting a prime-to-three Galois cover with its deck conjugate
gives a prime-to-three Galois cover of $C$. Its abelianization
would make $[Q'-P']$ three-primary by B7, contradicting both
possibilities in (B14.5). All constant extensions used here are
finite, and equality of sections restricts along their actual
Galois subgroups.
Finally, conjugate-commensurate decomposition groups would become
equal after restricting to the open subgroup of the constant-field
Galois group which is the image of their intersection. Their point
sections would then be equal over a finite extension, which we
have excluded. This proves every assertion of B14.
\end{proof}
